\chapter{Mathematical Background}
\label{ch:MathBackground}

This Chapter introduces the mathematical tools used in the Thesis: mixed quantum states, Fisher information, the symmetric logarithmic derivative, fidelity-based QFI bounds, and the quantum natural gradient. Supplementary material on generalized measurements, geometric extensions, and multiparameter bounds is collected in Appendix~\ref{app:mathematical_background}.

\section{Quantum States, Mixedness, and Open Systems}

\subsection{Density Operators, Purity, and Ontic versus Epistemic Uncertainty}

The following definitions establish how pure and mixed states are represented, distinguish intrinsic from classical uncertainty, and quantify mixedness for the reduced-state and QFI analyses used later in the Thesis.

\begin{definition}[Density operator and purity]
    A quantum state is represented by a density operator $\rho$ on a Hilbert space $\mathcal{H}$~\cite{nielsen_chuang_2010_quantum_computation_and_quantum_information, Sakurai_Napolitano_2020}:
    \begin{equation}
        \rho^\dagger = \rho \qquad \rho \geq 0 \qquad \operatorname{Tr}(\rho)=1
    \end{equation}
    Its spectral decomposition is
    \begin{equation}
        \rho = \sum_k p_k |\psi_k\rangle\langle\psi_k| \qquad p_k\geq0\qquad \sum_k p_k=1
    \end{equation}
    and its purity is
    \begin{equation}
        \gamma(\rho) \coloneqq \operatorname{Tr}(\rho^2)=\sum_kp_k^2
    \end{equation}
\end{definition}

\begin{definition}[Pure and mixed states]
    A pure state has the following properties:
    \begin{equation}
        \rho = |\psi_k\rangle\langle\psi_k|
        \qquad \gamma=1
    \end{equation}
    A mixed state has $\gamma<1$ and represents a probabilistic mixture of pure states.
\end{definition}

\begin{remark}[Ontic versus Epistemic Uncertainty]
    The distinction between pure and mixed states highlights two fundamentally different types of uncertainty. A pure state represents a perfectly known system. Any unpredictability in measurement is purely intrinsic to quantum mechanics (it is an \emph{ontic} uncertainty, from greek \emph{ontos}, meaning \enquote{being}). A mixed state introduces also classical uncertainty (\emph{epistemic} uncertainty, from greek \emph{episteme}, meaning \enquote{knowledge}), meaning the exact state of the system is unknown. This lack of knowledge typically arises either because the system was randomly prepared in one of several possible pure states, or because the system is entangled with an unobservable environment. Because completely different preparation methods can result in the exact same mixed state, quantities like the quantum Fisher information depend solely on the final parameterized density operator $\rho(\theta)$, rather than the specific ensemble of pure states used to create it. Here, $\theta$ denotes the parameter of interest; more generally, a model may depend on a vector of parameters $\boldsymbol{\theta}$.
\end{remark}

\begin{definition}[Partial trace]
    \label{def:partial_trace}
    Let a composite system have Hilbert space $\mathcal{H}_S\otimes\mathcal{H}_E$, and let $\mathcal{L}(\mathcal{H})$ denote the space of linear operators on $\mathcal{H}$. Following Ref.~\cite{nielsen_chuang_2010_quantum_computation_and_quantum_information}, the partial trace over $E$ is the linear map $\operatorname{Tr}_E:\mathcal{L}(\mathcal{H}_S\otimes\mathcal{H}_E)\to\mathcal{L}(\mathcal{H}_S)$ characterized on product operators by
    \begin{equation}
        \operatorname{Tr}_E[A_S\otimes B_E]
        \coloneqq
        A_S\operatorname{Tr}[B_E].
        \label{eq:partial_trace_definition}
    \end{equation}
    Equivalently, for rank-one product operators,
    \begin{equation}
        \operatorname{Tr}_E
        \left[
            |s_1\rangle\langle s_2|\otimes|e_1\rangle\langle e_2|
            \right]
        =
        \langle e_2|e_1\rangle
        |s_1\rangle\langle s_2|.
    \end{equation}
    Linearity determines its action on every joint operator. For a joint state $\rho_{SE}$, the reduced state available on $S$ is $\rho_S\coloneqq\operatorname{Tr}_E[\rho_{SE}]$.
\end{definition}

\begin{remark}[Operational meaning of the partial trace]
    Tracing out $E$ discards information available only through that subsystem while preserving all measurement statistics accessible on $S$. Thus, for every observable $A_S$ on $S$, $\operatorname{Tr}[\rho_S A_S]=\operatorname{Tr}[\rho_{SE}(A_S\otimes\mathbb{I}_E)]$. Even when $\rho_{SE}$ is pure, the reduced state $\rho_S$ can be mixed if $S$ and $E$ are entangled.
\end{remark}

\subsection{Parameterized States and Measurements}

\begin{definition}[Parameterized quantum model]
    Let $\rho_0$ be the initial state of the quantum system, and let $\theta\in\Theta\subseteq\mathbb{R}$ denote the unknown scalar parameter of interest. The parameter is encoded in the state through a parameter-dependent quantum channel $\mathcal{E}_\theta$, giving
    \begin{equation}
        \rho(\theta)
        = \mathcal{E}_\theta(\rho_0)
        = \sum_n K_n(\theta)\rho_0K_n^\dagger(\theta)
    \end{equation}
    where $\{K_n(\theta)\}_n$ is a set of $n$ Kraus operators. For a trace-preserving channel, they satisfy $\sum_n K_n^\dagger(\theta)K_n(\theta)=\mathbb{I}$.
    For a multiparameter model, one instead writes $\boldsymbol{\theta}\in\Theta\subseteq\mathbb{R}^p$ and $\rho(\boldsymbol{\theta})$.
\end{definition}


\begin{definition}[Positive operator-valued measure]
    Let $\mathcal{X}$ be the discrete set of possible measurement outcomes. A quantum measurement on $\mathcal{H}$ is described by a positive operator-valued measure (POVM), namely a family of operators $\{M_x\}_{x\in\mathcal{X}}$ satisfying
    \begin{equation}
        M_x\geq 0 \quad \text{for every }x\in\mathcal{X}
        \qquad
        \sum_{x\in\mathcal{X}} M_x=\mathbb{I}_{\mathcal{H}}.
    \end{equation}
    Here, $x$ is an outcome label, such as a detector port or a measured bit string. When the state is $\rho(\theta)$, the probability of observing $x\in\mathcal{X}$ is
    \begin{equation}
        p(x|\theta)
        =\operatorname{Tr}\!\left[\rho(\theta)M_x\right]
        \label{eq:povm_def}
    \end{equation}
    For a continuous outcome space, a POVM instead assigns a positive operator $M(B)$ to each measurable region $B\subseteq\mathcal{X}$, with $M(\mathcal{X})=\mathbb{I}_{\mathcal{H}}$ and $\Pr(x\in B\mid\theta)=\operatorname{Tr}[\rho(\theta)M(B)]$. Projective measurements are a special case of POVMs. Further details on projective measurements and general POVMs are provided in Appendix~\ref{app:mathematical_background}.
\end{definition}

The two definitions above give the basic quantum-to-classical sequence used throughout the Thesis.  The first arrow prepares a parameterized quantum state.  The measurement then converts that state into a classical probability distribution.
\begin{equation}
    \underbrace{
        \rho_0
        \xrightarrow{\mathcal{E}_\theta}
        \rho(\theta)
    }_{\text{quantum state and evolution}}
    \xrightarrow{\text{POVM }\{M_x\}_{x\in\mathcal{X}}}
    \underbrace{
        p(x\mid\theta)
    }_{\text{classical outcome probabilities}}
    \label{eq:quantum_to_classical_sequence}
\end{equation}
Noise and restricted access that act before the POVM belong inside the quantum part of this sequence.  Detector processing or statistical noise that acts on $p(x|\theta)$ belongs to the classical part.

\section{From Classical to Quantum Fisher Information}

This section links classical and quantum Fisher information: CFI quantifies the sensitivity of an observed probability distribution to changes in an unknown parameter, while QFI gives the greatest such sensitivity attainable over all quantum measurements.

\begin{definition}[Classical Fisher Information]
    Let $p(x|\theta)$ be the probability density of the measurement outcome $x$, with outcome space $\mathcal{X}$. The classical Fisher information (CFI) associated with this distribution is
    \begin{equation}
        \mathcal{F}_{\mathrm{C}}(\theta)
        \coloneqq
        \int_{\mathcal{X}}
        \frac{\left[\partial_\theta p(x|\theta)\right]^2}
        {p(x|\theta)}
        \,\mathrm{d}x
        \label{eq:classical_fi}
    \end{equation}
    For discrete outcomes, the integral is replaced by a sum over $x$.
\end{definition}

\begin{remark}[Meaning of CFI]
    Suppose that the measurement is repeated independently $N$ times. The variance of any locally unbiased estimator $\hat{\theta}$ satisfies the classical Cramér--Rao bound
    \begin{equation}
        \operatorname{Var}(\hat{\theta})
        \geq
        \frac{1}{N\mathcal{F}_{\mathrm{C}}(\theta)}
    \end{equation}
    So, the CFI quantifies the local sensitivity of the measurement statistics to changes in $\theta$: a larger CFI corresponds to a potentially smaller estimation variance. Because $p(x|\theta)$ is determined by the chosen POVM, the CFI is measurement dependent.
\end{remark}

Quantum Fisher information (QFI), unlike the CFI, does not depend on a predetermined measurement, it is intrinsic to the quantum state.

\begin{definition}[Quantum Fisher Information]
    For a differentiable, one-parameter family of quantum states $\rho(\theta)$, the Quantum Fisher Information (QFI) is obtained by maximizing the CFI over all allowed POVMs $\{M_x\}$:
    \begin{equation}
        \mathcal{F}_{\mathrm{Q}}(\theta)
        \coloneqq
        \sup_{\{M_x\}}
        \mathcal{F}_{\mathrm{C}}
        \bigl(\theta;{M_x}\bigr)
        \label{eq:qfi_def}
    \end{equation}
    \faLightbulb \hspace{5pt} The QFI provides a measurement-independent upper bound on the information that can be extracted locally from the state family $\rho(\theta)$.
\end{definition}

\begin{remark}[Quantum Cramér--Rao bound]
    Def. \ref{eq:qfi_def} is also known as Braunstein–Caves optimization.
    This result implies a quantum version of the Cramér--Rao bound (QCRB):
    \begin{equation}
        \operatorname{Var}(\hat{\theta})
        \geq\frac{1}{N\mathcal{F}_Q(\theta)}
        \label{eq:qcrb}
    \end{equation}
    The QCRB is the ideal precision benchmark. Practical receivers and variational bounds need to stay below this limit.
\end{remark}

\section{Mixed-State Quantum Fisher Information}

\subsection{Symmetric Logarithmic Derivative}

The optimization over all POVMs in Definition~\ref{eq:qfi_def} can be evaluated directly from the density operator by introducing the symmetric logarithmic derivative\footnote{It also explains why we defined QFI on a differentiable family of quantum states in Eq.~\eqref{eq:qfi_def}.}.

\begin{definition}[Symmetric logarithmic derivative]
    For a differentiable quantum state family $\rho(\theta)$, the symmetric logarithmic derivative (SLD) $L_\theta$ is an Hermitian operator that satisfies
    \begin{equation}
        \partial_\theta\rho(\theta)
        =\frac{1}{2}\{\rho(\theta),L_\theta\}
        =\frac{1}{2}\left[\rho(\theta)L_\theta+L_\theta\rho(\theta)\right]
        \label{eq:sld_lyapunov}
    \end{equation}
    where $\{A,B\}\coloneqq AB+BA$ is the anticommutator. For a rank-deficient state, the components of $L_\theta$ acting entirely within the kernel of $\rho$ are irrelevant and may be chosen arbitrarily.
\end{definition}

\begin{remark}[Intuitive meaning of the SLD]
    The SLD is the operator analogue of the classical score $\partial_\theta\log p(x\mid\theta)$: it describes how sensitively the quantum state changes with $\theta$. The symmetric product is required because $\rho$ and $L_\theta$ need not commute; when they do commute, the relation reduces to $\partial_\theta\rho=\rho L_\theta$, making the analogy with a logarithmic derivative explicit. Taking the trace gives $\operatorname{Tr}[\rho L_\theta]=0$, so the QFI in the following proposition is the second moment, or variance, of this quantum score.
\end{remark}

\begin{proposition}[QFI in terms of the Symmetric logarithmic derivative]
    For a differentiable one-parameter quantum model, the QFI can be evaluated as
    \begin{equation}
        \mathcal{F}_{\mathrm{Q}}(\theta)
        =
        \operatorname{Tr}
        \left[
            \rho(\theta)L_\theta^2
            \right]
        \label{eq:qfi_sld}
    \end{equation}
\end{proposition}

\begin{definition}[Quantum Fisher information matrix]
    For a model depending on a parameter vector
    $\boldsymbol{\theta}=(\theta_1,\ldots,\theta_m)$, the Quantum Fisher Information Matrix (QFIM) $H_{ij}$is defined by
    \begin{equation}
        H_{ij}(\boldsymbol{\theta})
        \coloneqq
        \frac{1}{2}
        \operatorname{Tr}
        \left[
            \rho(\boldsymbol{\theta})
            \left(L_iL_j+L_jL_i\right)
            \right]
        \label{eq:qfim}
    \end{equation}

    with $L_i$ denoting the Symmetric logarithmic derivative associated with $\theta_i$.

    For a one-parameter model, the QFI matrix reduces to the scalar QFI:
    \begin{equation}
        \mathcal{F}_{\mathrm{Q}}(\theta)
        =
        H_{\theta\theta}
    \end{equation}
\end{definition}

\begin{remark}[Intuitive meaning of the QFIM elements]
    The QFI is a sort of geometric map of the parameter space.
    \begin{itemize}
        \item \textbf{Diagonal elements ($H_{ii}$):} These represent the scalar QFI for individual parameters. They quantify the absolute sensitivity of the quantum state to variations in $\theta_i$ alone, assuming all other parameters are perfectly known.
        \item \textbf{Off-diagonal elements ($H_{ij}$ for $i \neq j$):} These capture the statistical correlation (cross-sensitivity), between the estimation of $\theta_i$ and $\theta_j$. If we have a non-zero off-diagonal term, it means that the estimators for these two parameters are statistically coupled. This coupling can introduce trade-offs, where optimizing the measurement for one parameter lowers the precision attainable for the other.
    \end{itemize}
\end{remark}

\subsection{Spectral and Vectorized Evaluation of QFI}

\begin{remark}[Exact QFI evaluation]
    Let the spectral decomposition of the density operator be
    \begin{equation}
        \rho
        =
        \sum_k
        \lambda_k
        |\psi_k\rangle\langle\psi_k|
    \end{equation}
    In the eigenbasis of $\rho$, the matrix elements of the Symmetric logarithmic derivative are
    \begin{equation}
        \langle\psi_k|L_\theta|\psi_l\rangle
        =\frac{2\langle\psi_k|\partial_\theta\rho|\psi_l\rangle}
        {\lambda_k+\lambda_l}
        \qquad \lambda_k+\lambda_l>0
        \label{eq:sld_spectral_elements}
    \end{equation}
    which yields
    \begin{equation}
        \mathcal{F}_{\mathrm{Q}}(\theta)
        =2\sum_{k,l:\lambda_k+\lambda_l>0}
        \frac{|\langle\psi_k|\partial_\theta\rho|\psi_l\rangle|^2}
        {\lambda_k+\lambda_l}
        \label{eq:qfi_spectral_explicit}
    \end{equation}
    This expression is exact but requires diagonalizing a $d\times d$ density matrix, with cost $\mathcal{O}(d^3)$.
\end{remark}
\begin{remark}[Vectorized evaluation of the QFIM]
    For the small systems used to obtain classical reference values, the QFIM may also be evaluated through the support-restricted vectorized formula~\cite{safranek_2018}
    \begin{equation}
        H_{ij}
        =2\,\operatorname{vec}[\partial_i\rho]^\dagger
        \left(\rho^T\otimes\mathbb{I}+\mathbb{I}\otimes\rho\right)^+
        \operatorname{vec}[\partial_j\rho]
        \label{eq:safranek_vectorized}
    \end{equation}
    Here, $\operatorname{vec}$ stacks the columns of a matrix into one column vector. With this convention,
    \begin{equation}
        \operatorname{vec}(ABC)
        =
        \left(C^T\otimes A\right)\operatorname{vec}(B)
    \end{equation}
    The two Kronecker-product terms therefore represent multiplication by $\rho$ from the right and from the left. The symbol ${}^{+}$ denotes the Moore--Penrose pseudoinverse. It avoids inverting the zero-eigenvalue sector of the density operator. This is the formula used for the model QFI calculation in Chapter~\ref{ch:exoplanet_experiment}.
\end{remark}

\section{Fidelity and the Bures Geometry}
Fidelity is a similarity measure between two quantum states. A fidelity close to $1$ indicates that the states are very similar and therefore difficult to distinguish, whereas a smaller fidelity indicates that they are more distinguishable. To make this intuitive notion of similarity precise for arbitrary pure or mixed quantum states, we use the root Uhlmann fidelity.

The Bures distance converts this measure of similarity into a geometric distance. Consequently, states with high fidelity are close in Bures distance, while states with low fidelity are farther apart and easier to distinguish.

\begin{definition}[Root fidelity and Bures distance]
    Let $\rho$ and $\sigma$ be density operators on the same Hilbert space. Their root Uhlmann fidelity is defined as
    \begin{equation}
        F(\rho,\sigma)\coloneqq
        \operatorname{Tr}\!\left[\sqrt{\sqrt{\rho}\sigma\sqrt{\rho}}\right]
        \label{eq:root_fidelity}
    \end{equation}
    In this Thesis, $F$ denotes the unsquared root fidelity. The corresponding Bures distance is defined by
    \begin{equation}
        d_{\mathrm B}^2(\rho,\sigma)\coloneqq2-2F(\rho,\sigma)
        \label{def:bures_distance}
    \end{equation}
\end{definition}

\begin{definition}[Bures metric]
    Let $\rho(\boldsymbol{\theta})$ be a differentiable quantum-state family with $\boldsymbol{\theta}\in\Theta\subseteq\mathbb{R}^p$. The components $g_{ij}^{\mathrm B}(\boldsymbol{\theta})$ of the Bures metric are defined by the second-order expansion of the squared Bures distance between neighboring states:
    \begin{equation}
        d_{\mathrm B}^2
        \left(
        \rho(\boldsymbol{\theta}),
        \rho(\boldsymbol{\theta}+\mathrm{d}\boldsymbol{\theta})
        \right)
        =
        \sum_{i,j=1}^{p}
        g_{ij}^{\mathrm B}(\boldsymbol{\theta})
        \,\mathrm{d}\theta_i\,\mathrm{d}\theta_j
        +
        o\!\left(
        \lVert\mathrm{d}\boldsymbol{\theta}\rVert^2
        \right)
        \label{eq:bures_local_expansion}
    \end{equation}
\end{definition}

\begin{proposition}[Relation between the Bures metric and the QFIM]
    At regular points where the rank of $\rho(\boldsymbol{\theta})$ is locally constant, the QFIM $H_{ij}$ and the Bures metric satisfy
    \begin{equation}
        H_{ij}(\boldsymbol{\theta})
        =4g_{ij}^{\mathrm B}(\boldsymbol{\theta})
        \label{def:qfim_bures_relation}
    \end{equation}
    Thus, the QFIM measures how rapidly neighboring quantum states separate in Bures geometry as the parameters change.
\end{proposition}


\section{Ansatz Geometry and the Quantum Natural Gradient}

Quantum Fisher information appears in this Thesis in two conceptually
distinct settings. In quantum sensing, the parameters
$\boldsymbol{\theta}$ represent unknown physical quantities encoded in a
state $\rho(\boldsymbol{\theta})$. The corresponding QFI or QFIM
quantifies how accurately these quantities can be estimated.

Consider now a variational quantum circuit: a sequence of parameterized
quantum gates that prepares a family of states
$\rho_{\boldsymbol{\alpha}}$. Its parameters
$\boldsymbol{\alpha}=(\alpha_1,\ldots,\alpha_q)
    \in\mathcal{A}\subseteq\mathbb{R}^q$ are real, adjustable control
variables, typically gate-rotation angles, chosen by an optimization
algorithm. The associated QFIM describes the geometry of the family of
states that the circuit can prepare.

\begin{definition}[Ansatz quantum Fisher information matrix]
    Let $\rho_{\boldsymbol{\alpha}}$ be the state prepared by an ansatz, with circuit parameters $\boldsymbol{\alpha}$. Its \emph{ansatz QFIM} is defined by the SLDs $\ell_\mu$ satisfying $\partial_\mu\rho_{\boldsymbol{\alpha}}=\{\rho_{\boldsymbol{\alpha}},\ell_\mu\}/2$:
    \begin{equation}
        F^{\mathrm{ans}}_{\mu\nu}
        \coloneqq
        \frac{1}{2}
        \operatorname{Tr}
        \left[
            \rho_{\boldsymbol{\alpha}}
            \left(
            \ell_\mu\ell_\nu
            +
            \ell_\nu\ell_\mu
            \right)
            \right]
        \label{eq:ansatz_qfim}
    \end{equation}
    This matrix quantifies the local distinguishability of ansatz states produced by infinitesimal changes in the circuit parameters.

    For a pure ansatz $\rho_{\boldsymbol{\alpha}}=|\psi(\boldsymbol{\alpha})\rangle\langle\psi(\boldsymbol{\alpha})|$, $F^{\mathrm{ans}}_{\mu\nu}$ reduces to four times the Fubini--Study metric,
    \begin{equation}
        F^{\mathrm{ans}}_{\mu\nu}=4g^{\mathrm{FS}}_{\mu\nu}
        \qquad
        g^{\mathrm{FS}}_{\mu\nu}
        =\operatorname{Re}\!\left[
            \langle\partial_\mu\psi|\partial_\nu\psi\rangle
            -\langle\partial_\mu\psi|\psi\rangle
            \langle\psi|\partial_\nu\psi\rangle\right]
        \label{eq:fubini_study_def}
    \end{equation}
    The subtraction of the second term removes changes associated only with the global phase of the state, which has no physical significance.

    Although the ansatz QFIM and the physical QFIM have the same mathematical form, they play different roles. The matrix $F^{\mathrm{ans}}_{\mu\nu}$ defines a metric on the circuit-parameter space and is used to guide optimization. By contrast, $H_{ij}(\boldsymbol{\theta})$ and $\mathcal{F}_{\mathrm Q}(\theta)$ quantify information about unknown physical parameters encoded in a sensor state.
\end{definition}

\begin{remark}[Quantum natural gradient]
    A variational algorithm assigns a scalar objective function to each set of circuit parameters in order to quantify the quality of the state prepared by the circuit. When the objective is written so that smaller values correspond to better solutions, it is called a loss function and is denoted by $\mathcal{L}(\boldsymbol{\alpha})$. The optimizer seeks parameters that minimize this loss.

    Ordinary gradient descent updates the circuit parameters in the direction of decreasing loss:
    \begin{equation}
        \boldsymbol{\alpha}_{t+1}
        =
        \boldsymbol{\alpha}_t
        -
        \eta
        \nabla_{\boldsymbol{\alpha}}
        \mathcal{L}(\boldsymbol{\alpha}_t)
    \end{equation}
    where $\eta$ is the learning rate. This update treats the circuit parameters as coordinates in a Euclidean space and does not account for how strongly different parameter variations change the prepared quantum state.

    Quantum natural-gradient descent incorporates the local geometry of the ansatz-state family by preconditioning the ordinary gradient with the ansatz QFIM:
    \begin{equation}
        \boldsymbol{\alpha}_{t+1}
        =
        \boldsymbol{\alpha}_t
        -
        \eta
        \left[
            F^{\mathrm{ans}}(\boldsymbol{\alpha}_t)
            +
            \lambda\mathbb{I}
            \right]^{+}
        \nabla_{\boldsymbol{\alpha}}
        \mathcal{L}(\boldsymbol{\alpha}_t),
        \qquad
        \lambda\geq0
        \label{eq:qng_update}
    \end{equation}
    Here, $\lambda$ is a regularization parameter and ${}^{+}$ denotes the Moore--Penrose pseudoinverse. The regularization term stabilizes the update when the ansatz QFIM is singular or poorly conditioned, as may occur when some circuit parameters are redundant or have little effect on the prepared state.

    The quantum natural gradient rescales the ordinary gradient according to the local distinguishability of neighboring ansatz states. It is a geometry-aware optimization direction, an example can be found in Info Box~\ref{sec:vqse_qng_optimization}.
\end{remark}

\section{Connection to the Variational Algorithms}
The mathematical tools introduced in this Chapter provide the foundation for the variational estimation methods developed in Chapter~\ref{ch:methods}.
For an $n$-qubit mixed state, the density matrix has dimension $d=2^n$, and its exact spectral decomposition has a dense computational cost of order $\mathcal{O}(d^3)=\mathcal{O}(8^n)$. Exact mixed-state QFI evaluation therefore becomes classically intractable as the system size increases.

To avoid this bottleneck, Chapter~\ref{ch:methods} introduces two complementary families of fidelity-based QFI bounds. The Truncated Quantum Fisher Information (TQFI) bounds retain the dominant eigenspace of the state, while the Sub- and Super-Fidelity QFI (SSQFI) bounds use trace-based quantities that do not require explicit eigenspace truncation. Both approaches provide finite-difference bounds on the QFI without requiring direct evaluation of the SLD for the full density matrix.

The Variational Quantum State Eigensolver (VQSE) supplies the most important eigenvalues and eigenvectors required to evaluate the TQFI bounds. These quantities are then incorporated into the Variational Quantum Fisher Information Estimation (VQFIE) procedure, which uses fidelity measurements to bound the QFI and, in the probe-optimization setting, to identify circuit parameters that improve the sensitivity of the sensor. The present Chapter provides the information-geometric foundation for these methods, while their circuits, optimization procedures, computational costs, and hardware requirements are developed in Chapter~\ref{ch:methods}.
